\documentclass[11pt]{article}
\usepackage[margin=1in]{geometry}
\usepackage{amsmath,amssymb,amsthm}
\usepackage[hidelinks]{hyperref}

\newtheorem{definition}{Definition}
\newtheorem{lemma}{Lemma}
\newtheorem{theorem}{Theorem}
\newcommand{\code}[1]{\texttt{#1}}

\title{Lea Overleaf Feature Playground}
\author{}
\date{}

\begin{document}
\maketitle

\noindent Paste this entire file into \code{main.tex} in a fresh Overleaf
project, with the Lea extension active and the local stack running.
It uses ordinary LaTeX packages; no \code{lea-tags.sty} is needed.
Open the Lean pane. There are \textbf{five tagged items}: one definition,
two consumers of a shared lemma, one intentionally faulty revision target,
and one optional longer target for manual pause testing.

\paragraph{Recommended order.}
Formalize \code{mix\_energy}, then \code{mix\_quarter\_drop}, then
\code{mix\_half\_bound}. Use individual buttons for the revision test.
The click-through recipes at the end give the source edits to make.
Keep all Lea labels unchanged throughout the experiments.

\section{One definition and one shared supporting result}

\begin{definition}[Centered energy]\label{def:mix-energy}
% lea: define label=mix_energy
% lea: context={Define mix_energy with explicit real arguments in the order c, x, y, and body (x - c)^2 + (y - c)^2. This is a computational definition, not a proposition.}
For $c,x,y\in\mathbb{R}$, define
\[
  E_c(x,y)=(x-c)^2+(y-c)^2.
\]
\end{definition}

% Intentionally UNTAGGED: this is supporting mathematics, not a sixth
% document target. The next two targets ask Lea to reserve it as a shared
% result. Do not add a formalize marker here and do not list this LaTeX
% lemma's label in uses: uses lists tagged document dependencies.
\begin{lemma}[Mixing energy identity]\label{lem:mixing-energy-step}
For every $c,t,x,y\in\mathbb{R}$,
\begin{align*}
  &E_c\bigl((1-t)x+ty,\;tx+(1-t)y\bigr)\\
  &\qquad=E_c(x,y)-2t(1-t)(x-y)^2.
\end{align*}
This identity has \emph{no restriction} on $t$.
\end{lemma}
\begin{proof}
Expand the two squares on the left using the definition of $E_c$.
Collect the coefficients of $x^2$, $xy$, $y^2$, $cx$, $cy$, and $c^2$.
Subtracting $E_c(x,y)$ leaves $-2t(1-t)(x-y)^2$.
\end{proof}

\begin{theorem}[First consumer: quarter mixing]\label{thm:mix-quarter-drop}
% lea: formalize label=mix_quarter_drop uses={mix_energy}
% lea: context={Follow the associated proof by applying the untagged Mixing energy identity. Search the project library, then Mathlib. If this exact supporting result is absent, call request_project_dependency with mode deferred and name mixing_energy_step, kind theorem, explicit binders (c t x y : Real), the unrestricted identity stated in the lemma above, and the import of the existing mix_energy definition. Deferred mode is intentional for this workflow test even though the algebra is short. Use the host-assigned module and declaration. Do not duplicate this general identity as a private helper or replace its application with a direct expansion of the target.}
For every $c,x,y\in\mathbb{R}$,
\[
 E_c\left(\frac{3x+y}{4},\frac{x+3y}{4}\right)
 = E_c(x,y)-\frac{3}{8}(x-y)^2.
\]
\end{theorem}
\begin{proof}
% lea: proof-for=mix_quarter_drop
Apply the Mixing energy identity with $t=1/4$.
The mixed coordinates become $(3x+y)/4$ and $(x+3y)/4$, and
$2(1/4)(1-1/4)=3/8$. Simplify these expressions to obtain the claim.
\end{proof}

\begin{theorem}[Second consumer: midpoint mixing]\label{thm:mix-half-bound}
% lea: formalize label=mix_half_bound uses={mix_energy}
% lea: context={Follow the associated proof by applying the same unrestricted Mixing energy identity. Search the project library, then Mathlib. Register reuse through request_project_dependency for mixing_energy_step with kind theorem, explicit binders (c t x y : Real), the exact unrestricted identity above, and the import of the existing mix_energy definition; use mode deferred if it is not already available. Use its host-assigned import and declaration. Do not create a differently specialized shared result or independently reprove the general identity.}
For every $c,x,y\in\mathbb{R}$,
\[
 E_c\left(\frac{x+y}{2},\frac{x+y}{2}\right)\le E_c(x,y).
\]
\end{theorem}
\begin{proof}
% lea: proof-for=mix_half_bound
Apply the same Mixing energy identity with $t=1/2$. Both coordinates
become $(x+y)/2$, so the left side equals
$E_c(x,y)-\tfrac12(x-y)^2$.
Since $(x-y)^2\ge0$, this is at most $E_c(x,y)$.
\end{proof}

\section{Pause, revise the source, and resume}

\noindent\textbf{The next statement and its proof are deliberately wrong.}
They omit the restriction on the mixing parameter. This is the test input
for author feedback and source-obstruction pauses. The recipes below
contain the corrected version; they are replacement instructions, not
hypotheses of the current theorem.

\begin{theorem}[Revision target, version 0]\label{thm:mix-revision-bound}
% lea: formalize label=mix_revision_bound uses={mix_energy}
% lea: context={Use the shared unrestricted mixing_energy_step identity and inspect the sign argument in the associated proof. Translate the current statement literally; replacement examples in the test recipes are not assumptions. If the current claim requires an additional hypothesis, publish a blocking Lea Status finding explaining it. Do not silently insert a restriction on t or change the conclusion.}
For every $c,t,x,y\in\mathbb{R}$,
\[
 E_c\bigl((1-t)x+ty,\;tx+(1-t)y\bigr)\le E_c(x,y).
\]
\end{theorem}
\begin{proof}
% lea: proof-for=mix_revision_bound
By the Mixing energy identity the left side is
$E_c(x,y)-2t(1-t)(x-y)^2$.
For every real $t$, the factor $t(1-t)$ is nonnegative.
Subtracting a nonnegative quantity proves the inequality.
\end{proof}

\section{Optional manual pause target}

\begin{theorem}[Repeated quarter mixing]\label{thm:mix-iterated-bound}
% lea: formalize label=mix_iterated_bound uses={mix_energy,mix_quarter_drop}
% lea: context={Use sequences x and y indexed by natural numbers, with the two recurrences as hypotheses for every index. Follow the induction proof using the checked mix_quarter_drop theorem. Keep the source statement and proof explicit; do not assume the desired bound.}
Let $c\in\mathbb{R}$ and let $x,y:\mathbb{N}\to\mathbb{R}$ satisfy,
for every $n\in\mathbb{N}$,
\[
 x(n+1)=\frac{3x(n)+y(n)}4,
 \qquad y(n+1)=\frac{x(n)+3y(n)}4.
\]
Then for every $n\in\mathbb{N}$,
\[
 E_c\bigl(x(n),y(n)\bigr)\le E_c\bigl(x(0),y(0)\bigr).
\]
\end{theorem}
\begin{proof}
% lea: proof-for=mix_iterated_bound
Induct on $n$. At $n=0$ the inequality is reflexive.
For the successor step, substitute the two recurrence equations into
$E_c(x(n+1),y(n+1))$ and apply the quarter-mixing theorem:
\[
 E_c(x(n+1),y(n+1))
 =E_c(x(n),y(n))-\frac38(x(n)-y(n))^2
 \le E_c(x(n),y(n)).
\]
Combine this with the induction hypothesis by transitivity.
\end{proof}

\section{Click-through recipes and expected evidence}

\subsection*{A. Generate a shared lemma, then reuse it}
\begin{enumerate}
 \item Formalize \code{mix\_energy} and wait for its checked definition.
 \item Formalize \code{mix\_quarter\_drop}. Its context explicitly asks for
   a deferred shared result, so the intended flow is: target waits,
   supporting worker proves the identity, target continues automatically.
 \item Expand \textbf{Shared mathematics}. Inspect the identity's full
   signature, canonical module/path, verification state, and attempt.
   The path should be under \code{Shared/}. The definition remains its
   own document module.
 \item Formalize \code{mix\_half\_bound}. Look for \textbf{Reused} on the
   same shared identity, with both consumers listed under \textbf{Used by}.
   Inspect the two generated target files: each should import and apply
   that same public lemma. Its attempt count should not increase.
\end{enumerate}
\noindent The contexts request this flow, but model choices are not a
deterministic test harness. A direct algebra proof that compiles has not
demonstrated sharing: inspect the shared record and actual imports.
The untagged supporting lemma should not increase the five document targets.
If the helper was already verified in this project, immediate reuse is expected.

\subsection*{B. Pause on a source issue, fix it, and Resume}
\begin{enumerate}
 \item In Lea settings, enable \textbf{Pause formalization when Lea needs
   author input} and save \emph{before} starting this run.
 \item Click the individual \textbf{Formalize} button on
   \code{mix\_revision\_bound}. Do not use \textbf{Formalize all} for this
   step: it disables automatic source pauses.
 \item Look for a blocking finding about the missing condition on $t$.
   A concrete counterexample to version 0 is $c=0$, $t=2$, $x=0$, $y=1$:
   the original energy is $1$, and the mixed energy is $5$.
   The intended stop is a source-obstruction pause. If the model instead
   completes a verified disproof, that is a different outcome; use recipe D
   for a manual pause.
 \item Once the run has settled as paused, edit the \emph{live theorem},
   replacing its opening sentence with:
\begin{verbatim}
For every $c,t,x,y\in\mathbb{R}$ with $0\le t\le1$,
\end{verbatim}
   Keep its displayed inequality, labels, and Lea context. In its live
   proof replace the false sign sentence with:
\begin{verbatim}
Since $0\le t\le1$, both $t$ and $1-t$ are nonnegative,
so $2t(1-t)(x-y)^2\ge0$.
\end{verbatim}
 \item Let Overleaf save the source, then click \textbf{Resume} on this
   target. Inspect Lea Status: it should reassess the earlier finding
   against the new source and prove the corrected claim. The saved session
   and earlier work should remain, with a linked continuation run.
\end{enumerate}
\noindent Resume requires a compatible saved checkpoint. If it reports an
incompatibility, keep the error visible when assessing the test;
\textbf{Re-formalize} is a separate fresh-start action.

\subsection*{C. Re-formalize a changed, completed target}
After recipe B succeeds, edit \code{mix\_revision\_bound} again.
Remove the interval condition and replace its opening sentence and
displayed conclusion with:
\begin{verbatim}
For every $c,t,x,y\in\mathbb{R}$,
\[
 E_c\bigl((1-t)x+ty,\;tx+(1-t)y\bigr)
 = E_c(x,y)-2t(1-t)(x-y)^2.
\]
\end{verbatim}
Replace its entire proof body, keeping the proof association comment, with:
\begin{verbatim}
% lea: proof-for=mix_revision_bound
Apply the shared Mixing energy identity directly.
\end{verbatim}
Also replace its two Lea metadata lines with:
\begin{verbatim}
% lea: formalize label=mix_revision_bound uses={mix_energy}
% lea: context={Apply the checked shared mixing_energy_step
% lea: identity directly for every real t, using its canonical import.}
\end{verbatim}
The corrected theorem is an equality valid for
all real $t$. Wait for the pane to show the source is out of date, then
click \textbf{Re-formalize}, possibly in the item's overflow menu.
Check that the new Lean statement is the unrestricted equality, that it
uses the existing shared identity, and that no second shared proof attempt
was created. This tests a new formalization against revised source.

\subsection*{D. Manual Pause and Resume, with a proof-only revision}
Formalize \code{mix\_iterated\_bound} after the quarter theorem is checked.
After a completed agent turn, click \textbf{Pause} while it is still running,
and wait for the pause to settle. Add this sentence to its live source proof,
just before ``Combine this with the induction hypothesis'':
\begin{verbatim}
The last inequality follows because $(x(n)-y(n))^2\ge0$.
\end{verbatim}
Let the source save, then \textbf{Resume}. Inspect the preserved work and
updated source assessment. The theorem itself is unchanged. If the proof
finishes before Pause takes effect, the completed result should remain;
use a fresh run if you want to try the timing again.

\end{document}
